\section{Zadanie 5}
Niech $G$ będzie skończonym grafem spójnym na co najmniej dwóch wierzchołkach. Rozważmy łańcuch Markowa, gdzie przestrzeń stanów $\Omega(G)$ jest rodziną wszystkich zbiorów niezależnych w $G$. Krok ze stanu $M_t$ do następnego stanu działa następująco. Losujemy jednostajnie i niezależnie: $uv \in E(G)$, $p \in \{1, 2, 3\}$.
\begin{itemize}
    \item Jeśli $p = 1$, to $M'_t = M_t \setminus \{u, v\}$.
    \item Jeśli $p = 2$, to $M'_t = M_t \setminus \{u\} \cup \{v\}$.
    \item Jeśli $p = 3$, to $M'_t = M_t \setminus \{v\} \cup \{u\}$.
\end{itemize}
Jeśli $M'_t$ jest zbiorem niezależnym w $G$, to $M_{t+1} = M'_t$, a w przeciwnym przypadku $M_{t+1} = M_t$. Udowodnij, że ten łańcuch jest nieprzywiedlny, nieokresowy i jego rozkład stacjonarny jest jednostajny.

\subsection*{Rozwiązanie}
Niech $N = |E(G)|$ oznacza liczbę krawędzi w grafie.

\subsection*{1. Nieprzywiedlność}

Łańcuch jest nieprzywiedlny, jeśli z każdego stanu $S$ można dojść do każdego innego stanu $T$ w skończonej liczbie kroków. Ponieważ wykażemy później, że łańcuch jest odwracalny, wystarczy pokazać, że z dowolnego zbioru niezależnego $S$ można dojść do zbioru pustego $\emptyset$.

Wybierzmy dowolny wierzchołek $u \in S$.
\begin{enumerate}
    \item Ponieważ $G$ jest spójny i ma co najmniej 2 wierzchołki, istnieje krawędź $uv \in E(G)$.
    \item Ponieważ $S$ jest zbiorem niezależnym i $u \in S$, to z definicji $v \notin S$.
\end{enumerate}

Wtedy weźmy krok łańcucha, w którym wylosowano krawędź $uv$ oraz $p=1$.
$M' = S \setminus \{u, v\}$. Skoro $v \notin S$, operacja ta redukuje się do $M' = S \setminus \{u\}$. Zbiór $M'$ (jako podzbiór zbioru niezależnego) jest nadal niezależny, więc przejście jest akceptowane z prawdopodobieństwem co najmniej $\frac{1}{3N}$.

Powtarzając tę procedurę dla każdego wierzchołka w $S$, redukujemy zbiór aż do osiągnięcia stanu $\emptyset$. Zatem graf przejść jest spójny.

\subsection*{2. Nieokresowość}

Aby wykazać nieokresowość, wystarczy wskazać stan, dla którego prawdopodobieństwo pozostania w nim w jednym kroku jest dodatnie ($P(S \to S) > 0$).

Niech $S$ będzie dowolnym zbiorem niezależnym. Rozważmy dowolną krawędź $uv$.
\begin{itemize}
    \item \textbf{Przypadek A:} $u \notin S$ i $v \notin S$.
    Jeśli wylosujemy $p=1$, to $M' = S \setminus \{u, v\} = S$. Stan nie ulega zmianie.
    \item \textbf{Przypadek B:} $u \in S$ (wtedy $v \notin S$).
    Jeśli wylosujemy $p=3$ (operacja $M_t \setminus \{v\} \cup \{u\}$), to ponieważ $v \notin S$ i $u \in S$, otrzymujemy $M' = S$. Stan nie ulega zmianie.
\end{itemize}
W obu przypadkach istnieje niezerowe prawdopodobieństwo pozostania w stanie $S$. Zatem łańcuch jest nieokresowy.

\subsection*{3. Jednostajny rozkład stacjonarny}

Aby udowodnić, że rozkład stacjonarny $\pi$ jest jednostajny, wystarczy pokazać, że macierz przejścia jest symetryczna, czyli dla dowolnych stanów $A, B$:
\[ P(A \to B) = P(B \to A) \]
Przejścia zmieniające stan są możliwe tylko między zbiorami różniącymi się w obrębie dwóch wierzchołków

\textbf{Przypadek 1: Zmiana liczby wierzchołków ($|A| = |B| + 1$)} \\
Niech $A = B \cup \{u\}$, gdzie $u \notin B$. Aby przejście było możliwe, $A$ musi być niezależny, więc żaden sąsiad $u$ nie należy do $B$.
\begin{itemize}
    \item \textbf{Przejście $A \to B$ (usunięcie $u$):} Wymaga wylosowania krawędzi $uv$ incydentnej z $u$ oraz $p=1$ ($M' = A \setminus \{u, v\}$). Ponieważ $v \notin A$, usuwamy tylko $u$. Prawdopodobieństwo wynosi:
    \[ P(A \to B) = \sum_{v \in N(u)} \frac{1}{3N} = \frac{\deg(u)}{3N} \]
    \item \textbf{Przejście $B \to A$ (dodanie $u$):} Wymaga wylosowania krawędzi $uv$ incydentnej z $u$. Musimy dodać $u$. W stanie $B$ nie ma $u$ ani $v$.
    Dla $p=3$ operacja to $B \setminus \{v\} \cup \{u\} = B \cup \{u\} = A$.
    Dla $p=2$ operacja to $B \setminus \{u\} \cup \{v\} = B \cup \{v\} \neq A$.
    Zatem interesuje nas tylko $p=3$:
    \[ P(B \to A) = \sum_{v \in N(u)} \frac{1}{3N} = \frac{\deg(u)}{3N} \]
\end{itemize}
Otrzymujemy $P(A \to B) = P(B \to A)$.

\textbf{Przypadek 2: Wymiana wierzchołka} \\
Niech $A$ zawiera $u$, $B$ zawiera $v$, gdzie $uv \in E(G)$, a $A \setminus \{u\} = B \setminus \{v\}$.
\begin{itemize}
    \item \textbf{Przejście $A \to B$:} Wymaga wylosowania konkretnej krawędzi $uv$ oraz $p=2$ ($A \setminus \{u\} \cup \{v\}$).
    \[ P(A \to B) = \frac{1}{N} \cdot \frac{1}{3} = \frac{1}{3N} \]
    \item \textbf{Przejście $B \to A$:} Wymaga wylosowania krawędzi $uv$ oraz $p=3$ ($B \setminus \{v\} \cup \{u\}$).
    \[ P(B \to A) = \frac{1}{N} \cdot \frac{1}{3} = \frac{1}{3N} \]
\end{itemize}
Tutaj również $P(A \to B) = P(B \to A)$.

\textbf{Wniosek:}
Macierz przejścia jest symetryczna, co z poprzedniego zestawu z łańcuchów Markowa wiemy, że daje nam rozkład stacjnonarny, jednostajny.